\chapter{Your First CNA Game}
\label{ch:first-game}

This chapter walks through the smallest complete CNA program --- the minimal XNA-style game
skeleton given in CNA's own README --- line by line, as a concrete anchor for the framework
concepts explained more fully in the chapters that follow.

\section{The whole program}

\begin{lstlisting}
#include <memory>

#include "Microsoft/Xna/Framework/Game.hpp"
#include "Microsoft/Xna/Framework/Color.hpp"
#include "Microsoft/Xna/Framework/Graphics/GraphicsDeviceManager.hpp"
#include "Microsoft/Xna/Framework/Graphics/SpriteBatch.hpp"
#include "Microsoft/Xna/Framework/Graphics/Texture2D.hpp"

using namespace Microsoft::Xna::Framework;
using namespace Microsoft::Xna::Framework::Graphics;

class MyGame final : public Game {
public:
    MyGame()
        : graphics_(this)
    {
    }

protected:
    void LoadContent() override
    {
        spriteBatch_ = std::make_unique<SpriteBatch>(getGraphicsDeviceProperty());
        logo_ = std::make_unique<Texture2D>("assets/logo.png", getGraphicsDeviceProperty());
    }

    void Update(GameTime& gameTime) override
    {
        (void)gameTime;
        // Update game state here.
    }

    void Draw(const GameTime& gameTime) override
    {
        (void)gameTime;

        auto& device = getGraphicsDeviceProperty();
        device.Clear(CornflowerBlue);

        spriteBatch_->Begin();
        spriteBatch_->Draw(*logo_, 100.0f, 80.0f);
        spriteBatch_->End();

        device.Present();
    }

private:
    GraphicsDeviceManager graphics_;
    std::unique_ptr<SpriteBatch> spriteBatch_;
    std::unique_ptr<Texture2D> logo_;
};

int main()
{
    MyGame game;
    game.Run();
    return 0;
}
\end{lstlisting}

\section{Reading it top to bottom}

\textbf{The includes are the API surface, and nothing else.} Every header included is a real
XNA~4.0 header, under the real XNA namespace path --- there is no CNA-specific include
required just to boot a game. This is the API-mirroring goal from
Chapter~\ref{ch:design-philosophy} made visible: a developer who has written XNA or FNA code
before can read this file without learning a single CNA-specific concept first.

\textbf{\texttt{MyGame} inherits \cnaclass{Game}, and owns a
\cnaclass{GraphicsDeviceManager} constructed against itself.} This is the standard XNA
ownership pattern: \cnaclass{GraphicsDeviceManager} is what actually creates and owns the
\cnaclass{GraphicsDevice}, negotiating with the platform window and the selected graphics
backend on the game's behalf; the game class itself never constructs a
\cnaclass{GraphicsDevice} directly. Chapter~\ref{ch:game-loop} covers
\cnaclass{GraphicsDeviceManager}'s role in the boot sequence in full, and
Chapter~\ref{ch:graphicsdevice} covers what the resulting \cnaclass{GraphicsDevice} exposes.

\textbf{\texttt{LoadContent()} is where GPU-backed resources are created, not the
constructor.} By the time \texttt{LoadContent()} runs, the graphics backend has already been
selected and initialized by \cnaclass{GraphicsDeviceManager}, so it is safe to construct a
\cnaclass{SpriteBatch} against \texttt{getGraphicsDeviceProperty()} and to load a
\cnaclass{Texture2D} from a file path. Constructing either of these in \texttt{MyGame}'s own
constructor would run before a graphics backend exists to own them. This ordering
constraint --- and CNA's raw-asset-plus-JSON-descriptor content model that makes
\texttt{"assets/logo.png"} a literal, directly loadable file path rather than a compiled
\texttt{.xnb} content-pipeline output --- is covered fully in
Chapter~\ref{ch:content-and-assets}.

\textbf{\texttt{Update(GameTime\&)} and \texttt{Draw(const GameTime\&)} are the two halves of
every frame, and note the constness difference.} \texttt{Update} takes a mutable
\texttt{GameTime\&} (some XNA game loops advance simulation-only clocks here);
\texttt{Draw} takes a \texttt{const GameTime\&}, since rendering should not mutate timing
state. Chapter~\ref{ch:game-loop} covers the full \cnaclass{Game} lifecycle this pair of
overrides plugs into --- \texttt{Initialize()}, \texttt{LoadContent()}, the
\texttt{Update} / \texttt{Draw} loop itself, and shutdown.

\textbf{The draw call is the canonical \cnaclass{SpriteBatch} triad:
\texttt{Begin()} / \texttt{Draw(...)} / \texttt{End()}.} \cnaclass{GraphicsDevice::Clear} takes a
named XNA color constant (\cnaclass{Color::CornflowerBlue} --- traditionally XNA's default
clear color in generated project templates, preserved here rather than swapped for something
``more CNA''), and \cnaclass{GraphicsDevice::Present()} is the final call of the frame,
regardless of which of the eleven backends from Part~IV is actually doing the presenting.
Chapter~\ref{ch:spritebatch} covers the full \cnaclass{SpriteBatch} API, including the
overloads of \texttt{Draw} beyond the two-float-position form used here.

\textbf{\texttt{main()} is exactly two lines of game-specific code.} Construct the game,
call \texttt{Run()}. Everything about window creation, backend initialization, the timing
loop, and shutdown lives inside \cnaclass{Game::Run()} --- covered in
Chapter~\ref{ch:game-loop} --- not in application code. This is deliberate: it is the same
shape a real XNA \texttt{Program.cs}'s \texttt{Main} method has, preserved so that porting an
existing XNA game's entry point is close to a mechanical translation.

\section{A second program: making it move}
\label{sec:first-game-movement}

The program above proves a texture can reach the screen; it says nothing about the
\texttt{Update} / \texttt{Draw} split actually \emph{doing} anything, since its own
\texttt{Update} is empty. This section extends exactly the same skeleton --- same includes,
same class shape, same \texttt{main()} --- with the smallest genuinely interactive addition:
moving the logo with the arrow keys, at a constant real-world speed regardless of frame rate.
Three small, real changes, each grounded in the actual headers rather than invented:

\begin{lstlisting}
#include "Microsoft/Xna/Framework/Input/Keyboard.hpp"
#include "Microsoft/Xna/Framework/Vector2.hpp"

using namespace Microsoft::Xna::Framework::Input;

class MyGame final : public Game {
    // ... constructor and LoadContent() unchanged ...

    void Update(GameTime& gameTime) override
    {
        const KeyboardState keys = Keyboard::GetState();
        const float seconds =
            static_cast<float>(gameTime.getElapsedGameTimeProperty().getTotalSecondsProperty());
        const float speed = 200.0f; // pixels per second

        if (keys.IsKeyDown(Keys::Left))  position_.X -= speed * seconds;
        if (keys.IsKeyDown(Keys::Right)) position_.X += speed * seconds;
        if (keys.IsKeyDown(Keys::Up))    position_.Y -= speed * seconds;
        if (keys.IsKeyDown(Keys::Down))  position_.Y += speed * seconds;
    }

    void Draw(const GameTime& gameTime) override
    {
        (void)gameTime;
        auto& device = getGraphicsDeviceProperty();
        device.Clear(CornflowerBlue);

        spriteBatch_->Begin();
        spriteBatch_->Draw(*logo_, position_, Color::White);
        spriteBatch_->End();

        device.Present();
    }

private:
    GraphicsDeviceManager graphics_;
    std::unique_ptr<SpriteBatch> spriteBatch_;
    std::unique_ptr<Texture2D> logo_;
    Vector2 position_{100.0f, 80.0f}; // new: replaces the two hardcoded floats
};
\end{lstlisting}

\textbf{Reading polled input is a static call, not an object you own.} \cnaclass{Keyboard} has
no constructor to call and nothing to store between frames --- \texttt{Keyboard::GetState()}
is a static method returning a fresh, immutable \cnaclass{KeyboardState} snapshot every time
it is called, and \texttt{Update} calls it once per frame. This is the same shape
Chapter~\ref{ch:input-system} covers for \cnaclass{Mouse} and \cnaclass{GamePad}.

\textbf{Frame-rate independence comes from \texttt{GameTime}, not a hardcoded per-frame
step.} Multiplying a pixels-per-second constant by
\texttt{gameTime.get\allowbreak{}ElapsedGameTimeProperty().get\allowbreak{}TotalSecondsProperty()} --- the real
\cnaclass{TimeSpan} accessor pair, not an invented shortcut --- means the logo moves at the
same real-world speed whether the game is running at 30 FPS or 240 FPS. A version that instead
moved by a flat constant every \texttt{Update} call would move four times faster on a display
with a four times higher refresh rate, which is precisely the class of bug frame-rate-independent
movement exists to avoid.

\begin{sourcenote}
\textbf{Why this program uses a different \texttt{Draw} overload than the first one.} The
first program in this chapter calls \texttt{spriteBatch\_->Draw(*logo\_, 100.0f, 80.0f)} ---
but that specific two-float-position overload, reading
\texttt{include/\allowbreak{}Microsoft/\allowbreak{}Xna/\allowbreak{}Framework/\allowbreak{}Graphics/\allowbreak{}SpriteBatch.hpp} directly, is tagged
\texttt{NOXNA}: a CNA-only convenience overload, not part of real XNA~4.0
(Chapter~\ref{ch:design-philosophy} covers what that tag means). This second program instead
calls \texttt{Draw(*logo\_, position\_, Color::White)}, the real, unmarked XNA~4.0 overload
that takes a \cnaclass{Vector2} position and an explicit tint color --- worth calling out
directly here, since a reader who wants code that ports cleanly to real XNA/FNA should reach
for this overload as the default, not the convenience one the first, minimal example used to
stay as short as possible.
\end{sourcenote}

\section{What happens if you change the backend}

Nothing in the file above mentions a graphics backend at all --- that is the entire point of
Chapter~\ref{ch:backend-architecture}'s abstraction layer. Recompiling this exact,
unmodified source file with \texttt{-DCNA\_GRAPHICS\allowbreak{}\_BACKEND=VULKAN} instead of
\texttt{EASYGL} produces a program that opens a Vulkan-backed window, uploads
\texttt{assets/logo.png} through the Vulkan texture path, and issues the same
\cnaclass{SpriteBatch} draw call through Vulkan command buffers instead of OpenGL calls ---
with the game-facing source unchanged. Part~IV of this volume is dedicated to what actually
differs underneath that unchanged surface, backend by backend.
